\documentclass[]{article}
\usepackage{xeCJK}
\setCJKmainfont{Noto Serif SC}
\usepackage{fontspec}
\usepackage{amsmath}
\usepackage{unicode-math}
\usepackage{graphicx}
\usepackage{geometry}
\geometry{}
\makeatletter
\providecommand{\tightlist}{\itemsep0pt}
\providecommand{\itemize}{\itemsep0pt}
\makeatother
\begin{document}
\section{文件样式踩坑测试}\label{ux6587ux4ef6ux6837ux5f0fux8e29ux5751ux6d4bux8bd5}

\begin{quote}
这是用来测试各种文件编辑场景的测试文件，涵盖 Obsidian
笔记中常见的所有元素。
\end{quote}

\begin{center}\rule{0.5\linewidth}{0.5pt}\end{center}

\subsection{1.
普通文本与中文}\label{ux666eux901aux6587ux672cux4e0eux4e2dux6587}

这是普通段落文字。中文段落测试：的概率与统计的交汇融合，构成了极具选拔功能的''28分高分值协同效应区''。

\subsubsection{1.1 加粗与斜体}\label{ux52a0ux7c97ux4e0eux659cux4f53}

\textbf{加粗文字}用于强调重点。\emph{斜体文字}用于次级强调。\textbf{\emph{加粗斜体}}组合使用。

中文强调测试：\textbf{高可靠性的深层传输网络}，需要在设计时考虑多个因素。

\begin{center}\rule{0.5\linewidth}{0.5pt}\end{center}

\subsection{2. 标题层级}\label{ux6807ux9898ux5c42ux7ea7}

\subsubsection{三级标题}\label{ux4e09ux7ea7ux6807ux9898}

\paragraph{四级标题}\label{ux56dbux7ea7ux6807ux9898}

\subparagraph{五级标题}\label{ux4e94ux7ea7ux6807ux9898}

六级标题

\begin{center}\rule{0.5\linewidth}{0.5pt}\end{center}

\subsection{3. 列表}\label{ux5217ux8868}

\subsubsection{3.1 无序列表}\label{ux65e0ux5e8fux5217ux8868}

\begin{itemize}
\tightlist
\item
  第一条要点

  \begin{itemize}
  \tightlist
  \item
    子要点（缩进两个空格）

    \begin{itemize}
    \tightlist
    \item
      更深的子要点
    \end{itemize}
  \end{itemize}
\item
  第二条要点

  \begin{itemize}
  \tightlist
  \item
    嵌套内容
  \end{itemize}
\end{itemize}

\subsubsection{3.2 有序列表}\label{ux6709ux5e8fux5217ux8868}

\begin{enumerate}
\def\labelenumi{\arabic{enumi}.}
\tightlist
\item
  第一步操作
\item
  第二步操作

  \begin{enumerate}
  \def\labelenumii{\arabic{enumii}.}
  \tightlist
  \item
    子步骤 A
  \item
    子步骤 B
  \end{enumerate}
\item
  第三步操作
\end{enumerate}

\subsubsection{3.3 任务列表（Obsidian
特性）}\label{ux4efbux52a1ux5217ux8868obsidian-ux7279ux6027}

\begin{itemize}
\tightlist
\item[$\square$]
  未完成的任务
\item[$\boxtimes$]
  已完成的任务（Obsidian 中显示为勾选框）
\item[$\square$]
  另一个待办
\end{itemize}

\begin{center}\rule{0.5\linewidth}{0.5pt}\end{center}

\subsection{4. 代码}\label{ux4ee3ux7801}

\subsubsection{4.1 行内代码}\label{ux884cux5185ux4ee3ux7801}

Python 中定义变量 \texttt{x\ =\ 100}，JavaScript 中
\texttt{let\ y\ =\ 200}。

\subsubsection{4.2
代码块（多语言）}\label{ux4ee3ux7801ux5757ux591aux8bedux8a00}

\begin{Shaded}
\begin{Highlighting}[]
\CommentTok{\# Python 代码块}
\KeywordTok{def}\NormalTok{ fibonacci(n):}
    \ControlFlowTok{if}\NormalTok{ n }\OperatorTok{\textless{}=} \DecValTok{1}\NormalTok{:}
        \ControlFlowTok{return}\NormalTok{ n}
    \ControlFlowTok{return}\NormalTok{ fibonacci(n}\OperatorTok{{-}}\DecValTok{1}\NormalTok{) }\OperatorTok{+}\NormalTok{ fibonacci(n}\OperatorTok{{-}}\DecValTok{2}\NormalTok{)}

\NormalTok{result }\OperatorTok{=}\NormalTok{ [fibonacci(i) }\ControlFlowTok{for}\NormalTok{ i }\KeywordTok{in} \BuiltInTok{range}\NormalTok{(}\DecValTok{10}\NormalTok{)]}
\BuiltInTok{print}\NormalTok{(result)}
\end{Highlighting}
\end{Shaded}

\begin{Shaded}
\begin{Highlighting}[]
\CommentTok{// JavaScript 代码块}
\KeywordTok{const}\NormalTok{ fibonacci }\OperatorTok{=}\NormalTok{ (n) }\KeywordTok{=\textgreater{}}\NormalTok{ \{}
    \ControlFlowTok{if}\NormalTok{ (n }\OperatorTok{\textless{}=} \DecValTok{1}\NormalTok{) }\ControlFlowTok{return}\NormalTok{ n}\OperatorTok{;}
    \ControlFlowTok{return} \FunctionTok{fibonacci}\NormalTok{(n}\OperatorTok{{-}}\DecValTok{1}\NormalTok{) }\OperatorTok{+} \FunctionTok{fibonacci}\NormalTok{(n}\OperatorTok{{-}}\DecValTok{2}\NormalTok{)}\OperatorTok{;}
\NormalTok{\}}\OperatorTok{;}

\BuiltInTok{console}\OperatorTok{.}\FunctionTok{log}\NormalTok{([}\OperatorTok{...}\BuiltInTok{Array}\NormalTok{(}\DecValTok{10}\NormalTok{)]}\OperatorTok{.}\FunctionTok{map}\NormalTok{((\_}\OperatorTok{,}\NormalTok{ i) }\KeywordTok{=\textgreater{}} \FunctionTok{fibonacci}\NormalTok{(i)))}\OperatorTok{;}
\end{Highlighting}
\end{Shaded}

\begin{Shaded}
\begin{Highlighting}[]
\CommentTok{\# Shell/Bash 代码块}
\BuiltInTok{cd} \StringTok{"C:\textbackslash{}Users\textbackslash{}39173\textbackslash{}Desktop\textbackslash{}笔记"}
\FunctionTok{git}\NormalTok{ status}
\FunctionTok{git}\NormalTok{ add .}
\FunctionTok{git}\NormalTok{ commit }\AttributeTok{{-}m} \StringTok{"update notes"}
\end{Highlighting}
\end{Shaded}

\begin{Shaded}
\begin{Highlighting}[]
\CommentTok{\% LaTeX 数学代码块}
\SpecialStringTok{\textbackslash{}[}
\SpecialCharTok{\textbackslash{}mathcal}\SpecialStringTok{\{L\} = }\SpecialCharTok{\textbackslash{}int}\SpecialStringTok{\_\{}\SpecialCharTok{\textbackslash{}Omega}\SpecialStringTok{\} |}\SpecialCharTok{\textbackslash{}nabla}\SpecialStringTok{ u|\^{}2 }\SpecialCharTok{\textbackslash{},}\SpecialStringTok{ dx}
\SpecialStringTok{\textbackslash{}]}
\end{Highlighting}
\end{Shaded}

\begin{center}\rule{0.5\linewidth}{0.5pt}\end{center}

\subsection{5. 数学公式（LaTeX）}\label{ux6570ux5b66ux516cux5f0flatex}

\subsubsection{5.1 行内公式}\label{ux884cux5185ux516cux5f0f}

当 \(p_1 = \frac{1}{3}\) 时，对任意 \(n \ge 1\)，均有
\(0 < p_n < \frac{1}{2}\)。

随机变量序列 \(z^{(1)}, z^{(2)}, \dots, z^{(M)}\) 满足马尔可夫性质。

\subsubsection{5.2 独立公式块}\label{ux72ecux7acbux516cux5f0fux5757}

条件概率：

\[P_n = \frac{2}{3} P_{n+1} + \frac{1}{3} P_{n-1} \quad (1 \le n \le 5)\]

递推公式推导：

\[E_n = \frac{2}{3} (E_{n+1} + 1) + \frac{1}{3} (E_{n-1} + 1) \quad (1 \le n \le 5)\]

全概率展开：

\[P(X = x) = \sum_{i=1}^{n} P(A_i) \cdot P(X = x \mid A_i)\]

极限分布：

\[\lim_{n \to \infty} p_n = \frac{1}{5}\]

分段函数：

\[
f(x) = 
\begin{cases}
\frac{1}{5} + \frac{2}{5}\lambda_1^n + \frac{2}{5}\lambda_2^n & \text{if } n \ge 0 \\
0 & \text{otherwise}
\end{cases}
\]

\begin{center}\rule{0.5\linewidth}{0.5pt}\end{center}

\subsection{6. 表格}\label{ux8868ux683c}

\subsubsection{6.1 普通表格}\label{ux666eux901aux8868ux683c}

\begin{longtable}[]{@{}llll@{}}
\toprule\noalign{}
知识模块 & 分值（分） & 占总分比例（\%） & 考查形式与命题动向 \\
\midrule\noalign{}
\endhead
\bottomrule\noalign{}
\endlastfoot
函数与导数 & 38 & 25.33 & 压轴解答题，侧重非线性方程根的分布 \\
解析几何 & 26 & 17.33 & 中高档解答题，计算负荷极高 \\
数列 & 17 & 11.33 & 重点考查递推通项结构性质 \\
概率与统计 & 11 & 7.33 & 侧重条件概率、全概率公式、期望 \\
\end{longtable}

\subsubsection{6.2 左对齐表格}\label{ux5de6ux5bf9ux9f50ux8868ux683c}

\begin{longtable}[]{@{}lll@{}}
\toprule\noalign{}
命题范式 & 核心模型 & 典型工具 \\
\midrule\noalign{}
\endhead
\bottomrule\noalign{}
\endlastfoot
一维随机游走 & 二阶常系数线性递推 & 累加法、待定系数法 \\
马尔可夫链 & 联立多元差分方程组 & 特征方程求特征根 \\
非线性递推 & 一阶非线性三次多项式 & 数学归纳法 \\
\end{longtable}

\begin{center}\rule{0.5\linewidth}{0.5pt}\end{center}

\subsection{7. 引用块}\label{ux5f15ux7528ux5757}

\begin{quote}
这是一段引用文字。

第二行引用。引用块常用于标注重要概念或引用来源。
\end{quote}

多层嵌套引用：

\begin{quote}
外层引用 \textgreater{} 内层引用 \textgreater\textgreater{} 更深的嵌套
\end{quote}

\begin{center}\rule{0.5\linewidth}{0.5pt}\end{center}

\subsection{8. 分隔线}\label{ux5206ux9694ux7ebf}

\begin{center}\rule{0.5\linewidth}{0.5pt}\end{center}

上面是分隔线，上面是分隔线。

\begin{center}\rule{0.5\linewidth}{0.5pt}\end{center}

\subsection{9. 链接}\label{ux94feux63a5}

\subsubsection{9.1 普通链接}\label{ux666eux901aux94feux63a5}

\href{https://portal.pku.edu.cn/portal2017/}{北京大学门户}

\href{https://ocw.mit.edu/}{MIT OpenCourseWare}

\subsubsection{9.2
维基风格链接}\label{ux7ef4ux57faux98ceux683cux94feux63a5}

{[}{[}常用网站{]}{]} --- 链接到同目录下的另一个笔记文件（Obsidian
内部链接）

\begin{center}\rule{0.5\linewidth}{0.5pt}\end{center}

\subsection{10. Obsidian 特有的 YAML
属性}\label{obsidian-ux7279ux6709ux7684-yaml-ux5c5eux6027}

\begin{Shaded}
\begin{Highlighting}[]
\PreprocessorTok{{-}{-}{-}}
\FunctionTok{uid}\KeywordTok{:}\AttributeTok{ 2026{-}05{-}21{-}test}
\FunctionTok{tags}\KeywordTok{:}
\AttributeTok{  }\KeywordTok{{-}}\AttributeTok{ 测试}
\AttributeTok{  }\KeywordTok{{-}}\AttributeTok{ 笔记}
\FunctionTok{created}\KeywordTok{:}\AttributeTok{ 2026{-}05{-}21}
\FunctionTok{modified}\KeywordTok{:}\AttributeTok{ 2026{-}05{-}21}
\FunctionTok{aliases}\KeywordTok{:}
\AttributeTok{  }\KeywordTok{{-}}\AttributeTok{ 文件样式测试}
\PreprocessorTok{{-}{-}{-}}

\AttributeTok{以上为 YAML frontmatter，Obsidian 会读取这些属性。}
\end{Highlighting}
\end{Shaded}

\begin{center}\rule{0.5\linewidth}{0.5pt}\end{center}

\subsection{11.
高亮与删除线}\label{ux9ad8ux4eaeux4e0eux5220ux9664ux7ebf}

这是普通文字 ==高亮文字== 用于突出重点。

这是普通文字 \st{删除线文字} 用于标记过时内容。

\begin{center}\rule{0.5\linewidth}{0.5pt}\end{center}

\subsection{12.
水平分割符的特殊情况}\label{ux6c34ux5e73ux5206ux5272ux7b26ux7684ux7279ux6b8aux60c5ux51b5}

单独一行 \texttt{-\/-\/-} 会渲染为分隔线。如果你想输出 \texttt{-\/-\/-}
文字而不转为分隔线，需要使用转义或在其他元素中。

\begin{center}\rule{0.5\linewidth}{0.5pt}\end{center}

\subsection{13.
转义字符测试}\label{ux8f6cux4e49ux5b57ux7b26ux6d4bux8bd5}

以下字符在 Markdown 中需要转义：

\begin{itemize}
\tightlist
\item
  星号：* （避免变为斜体）
\item
  下划线：\_ （避免变为斜体）
\item
  井号：\# （避免变为标题）
\item
  反引号：` （避免变为行内代码）
\end{itemize}

在 Obsidian 中，普通段落内单独一个 \texttt{-}
不会变列表，需要有空行或其他标记触发。

\begin{center}\rule{0.5\linewidth}{0.5pt}\end{center}

\subsection{14.
常见踩坑点总结}\label{ux5e38ux89c1ux8e29ux5751ux70b9ux603bux7ed3}

\begin{longtable}[]{@{}
  >{\raggedright\arraybackslash}p{(\linewidth - 4\tabcolsep) * \real{0.3333}}
  >{\raggedright\arraybackslash}p{(\linewidth - 4\tabcolsep) * \real{0.3333}}
  >{\raggedright\arraybackslash}p{(\linewidth - 4\tabcolsep) * \real{0.3333}}@{}}
\toprule\noalign{}
\begin{minipage}[b]{\linewidth}\raggedright
场景
\end{minipage} & \begin{minipage}[b]{\linewidth}\raggedright
坑
\end{minipage} & \begin{minipage}[b]{\linewidth}\raggedright
解决方式
\end{minipage} \\
\midrule\noalign{}
\endhead
\bottomrule\noalign{}
\endlastfoot
中文文件名 & 部分工具对中文路径处理不一致 & 始终用引号包裹路径 \\
表格内竖线 & \texttt{\textbar{}} 出现在表格内需要转义 &
表格内尽量不用竖线 \\
代码块中的 LaTeX & \texttt{\$\$} 可能与外层解析冲突 &
代码块内无需转义，直接写 \texttt{\$\$} \\
分隔线 \texttt{-\/-\/-} & 单独一行会变成横线 & 用其他方式表达，如
\texttt{-\/-\/-\ 文字\ -\/-\/-} \\
Obsidian 内部链接 & \texttt{{[}{[}笔记名{]}{]}} 在非 Obsidian 环境会失效
& 区分使用场景 \\
YAML frontmatter & 必须放在文件最开头，冒号后有空格 &
\texttt{key:\ value} 格式 \\
\end{longtable}

\begin{center}\rule{0.5\linewidth}{0.5pt}\end{center}

\subsection{15.
空文件与特殊文件名}\label{ux7a7aux6587ux4ef6ux4e0eux7279ux6b8aux6587ux4ef6ux540d}

测试一些特殊情况：

\begin{itemize}
\tightlist
\item
  文件名带空格：\texttt{未命名\ 1.md}
\item
  文件名带中文：\texttt{常用网站.md}
\item
  零字节文件：某些占位文件可能为空
\item
  Canvas 文件：\texttt{.canvas} 格式是 JSON 结构
\end{itemize}

\begin{center}\rule{0.5\linewidth}{0.5pt}\end{center}

\emph{测试文件 Created: 2026-05-21}
\emph{用于测试各种文件编辑场景与踩坑记录}
\end{document}